%% Gato pixel del sitio PAES (CAT_PX de paes.html), generado una vez.
%% \gato[escala]{cuerpo}{sombra}: colores de pelaje; \gatonaranjo, \gatogris, ... son atajos.
\definecolor{gatoK}{HTML}{1A1008}\definecolor{gatoP}{HTML}{FF9ECB}\definecolor{gatoY}{HTML}{FFD75E}\definecolor{gatoR}{HTML}{E8342A}
\newcommand{\gato}[3][0.07]{\begin{tikzpicture}[x=#1cm,y=#1cm,baseline=-8]
\colorlet{gatoB}{#2}\colorlet{gatob}{#3}
\fill[gatoK] (3,0) rectangle ++(2,-1);
\fill[gatoK] (15,0) rectangle ++(2,-1);
\fill[gatoK] (2,-1) rectangle ++(1,-1);
\fill[gatoB] (3,-1) rectangle ++(2,-1);
\fill[gatoK] (5,-1) rectangle ++(1,-1);
\fill[gatoK] (14,-1) rectangle ++(1,-1);
\fill[gatoB] (15,-1) rectangle ++(2,-1);
\fill[gatoK] (17,-1) rectangle ++(1,-1);
\fill[gatoK] (1,-2) rectangle ++(1,-1);
\fill[gatoB] (2,-2) rectangle ++(1,-1);
\fill[gatoP] (3,-2) rectangle ++(2,-1);
\fill[gatoB] (5,-2) rectangle ++(1,-1);
\fill[gatoK] (6,-2) rectangle ++(1,-1);
\fill[gatoK] (13,-2) rectangle ++(1,-1);
\fill[gatoB] (14,-2) rectangle ++(1,-1);
\fill[gatoP] (15,-2) rectangle ++(2,-1);
\fill[gatoB] (17,-2) rectangle ++(1,-1);
\fill[gatoK] (18,-2) rectangle ++(1,-1);
\fill[gatoK] (1,-3) rectangle ++(1,-1);
\fill[gatoB] (2,-3) rectangle ++(1,-1);
\fill[gatoP] (3,-3) rectangle ++(2,-1);
\fill[gatoB] (5,-3) rectangle ++(2,-1);
\fill[gatoK] (7,-3) rectangle ++(6,-1);
\fill[gatoB] (13,-3) rectangle ++(2,-1);
\fill[gatoP] (15,-3) rectangle ++(2,-1);
\fill[gatoB] (17,-3) rectangle ++(1,-1);
\fill[gatoK] (18,-3) rectangle ++(1,-1);
\fill[gatoK] (1,-4) rectangle ++(1,-1);
\fill[gatoB] (2,-4) rectangle ++(16,-1);
\fill[gatoK] (18,-4) rectangle ++(1,-1);
\fill[gatoK] (0,-5) rectangle ++(1,-1);
\fill[gatoB] (1,-5) rectangle ++(18,-1);
\fill[gatoK] (19,-5) rectangle ++(1,-1);
\fill[gatoK] (0,-6) rectangle ++(1,-1);
\fill[gatoB] (1,-6) rectangle ++(2,-1);
\fill[gatoK] (3,-6) rectangle ++(1,-1);
\fill[gatoK] (4,-6) rectangle ++(1,-1);
\fill[white] (5,-6) rectangle ++(1,-1);
\fill[gatoK] (6,-6) rectangle ++(1,-1);
\fill[gatoB] (7,-6) rectangle ++(6,-1);
\fill[gatoK] (13,-6) rectangle ++(1,-1);
\fill[gatoK] (14,-6) rectangle ++(1,-1);
\fill[white] (15,-6) rectangle ++(1,-1);
\fill[gatoK] (16,-6) rectangle ++(1,-1);
\fill[gatoB] (17,-6) rectangle ++(2,-1);
\fill[gatoK] (19,-6) rectangle ++(1,-1);
\fill[gatoK] (0,-7) rectangle ++(1,-1);
\fill[gatoB] (1,-7) rectangle ++(2,-1);
\fill[gatoK] (3,-7) rectangle ++(1,-1);
\fill[gatoK] (4,-7) rectangle ++(2,-1);
\fill[gatoK] (6,-7) rectangle ++(1,-1);
\fill[gatoB] (7,-7) rectangle ++(1,-1);
\fill[gatob] (8,-7) rectangle ++(1,-1);
\fill[gatoB] (9,-7) rectangle ++(2,-1);
\fill[gatob] (11,-7) rectangle ++(1,-1);
\fill[gatoB] (12,-7) rectangle ++(1,-1);
\fill[gatoK] (13,-7) rectangle ++(1,-1);
\fill[gatoK] (14,-7) rectangle ++(2,-1);
\fill[gatoK] (16,-7) rectangle ++(1,-1);
\fill[gatoB] (17,-7) rectangle ++(2,-1);
\fill[gatoK] (19,-7) rectangle ++(1,-1);
\fill[gatoK] (0,-8) rectangle ++(1,-1);
\fill[gatoB] (1,-8) rectangle ++(7,-1);
\fill[white] (8,-8) rectangle ++(4,-1);
\fill[gatoB] (12,-8) rectangle ++(7,-1);
\fill[gatoK] (19,-8) rectangle ++(1,-1);
\fill[gatoK] (0,-9) rectangle ++(1,-1);
\fill[gatoB] (1,-9) rectangle ++(1,-1);
\fill[gatob] (2,-9) rectangle ++(1,-1);
\fill[gatoB] (3,-9) rectangle ++(4,-1);
\fill[white] (7,-9) rectangle ++(2,-1);
\fill[gatoP] (9,-9) rectangle ++(2,-1);
\fill[white] (11,-9) rectangle ++(2,-1);
\fill[gatoB] (13,-9) rectangle ++(4,-1);
\fill[gatob] (17,-9) rectangle ++(1,-1);
\fill[gatoB] (18,-9) rectangle ++(1,-1);
\fill[gatoK] (19,-9) rectangle ++(1,-1);
\fill[gatoK] (1,-10) rectangle ++(1,-1);
\fill[gatoB] (2,-10) rectangle ++(6,-1);
\fill[white] (8,-10) rectangle ++(4,-1);
\fill[gatoB] (12,-10) rectangle ++(6,-1);
\fill[gatoK] (18,-10) rectangle ++(1,-1);
\fill[gatoK] (1,-11) rectangle ++(1,-1);
\fill[gatoB] (2,-11) rectangle ++(2,-1);
\fill[gatoY] (4,-11) rectangle ++(5,-1);
\fill[gatoR] (9,-11) rectangle ++(2,-1);
\fill[gatoY] (11,-11) rectangle ++(5,-1);
\fill[gatoB] (16,-11) rectangle ++(2,-1);
\fill[gatoK] (18,-11) rectangle ++(1,-1);
\fill[gatoK] (1,-12) rectangle ++(1,-1);
\fill[gatoB] (2,-12) rectangle ++(16,-1);
\fill[gatoK] (18,-12) rectangle ++(1,-1);
\fill[gatoK] (1,-13) rectangle ++(1,-1);
\fill[gatoB] (2,-13) rectangle ++(1,-1);
\fill[gatob] (3,-13) rectangle ++(1,-1);
\fill[gatoB] (4,-13) rectangle ++(12,-1);
\fill[gatob] (16,-13) rectangle ++(1,-1);
\fill[gatoB] (17,-13) rectangle ++(1,-1);
\fill[gatoK] (18,-13) rectangle ++(1,-1);
\fill[gatoK] (2,-14) rectangle ++(1,-1);
\fill[gatoB] (3,-14) rectangle ++(2,-1);
\fill[gatoK] (5,-14) rectangle ++(2,-1);
\fill[gatoB] (7,-14) rectangle ++(6,-1);
\fill[gatoK] (13,-14) rectangle ++(2,-1);
\fill[gatoB] (15,-14) rectangle ++(2,-1);
\fill[gatoK] (17,-14) rectangle ++(1,-1);
\fill[gatoK] (3,-15) rectangle ++(2,-1);
\fill[gatoK] (7,-15) rectangle ++(6,-1);
\fill[gatoK] (15,-15) rectangle ++(2,-1);
\end{tikzpicture}}
\definecolor{gNar}{HTML}{FF9A3C}\definecolor{gNarS}{HTML}{E07A1F}
\definecolor{gGri}{HTML}{A7A7AD}\definecolor{gGriS}{HTML}{7D7D84}
\definecolor{gBla}{HTML}{F4F1EA}\definecolor{gBlaS}{HTML}{D8D2C4}
\definecolor{gNeg}{HTML}{3A3340}\definecolor{gNegS}{HTML}{241F29}
\definecolor{gAti}{HTML}{C8935A}\definecolor{gAtiS}{HTML}{8A5A2B}
\definecolor{gLil}{HTML}{B57BE0}\definecolor{gLilS}{HTML}{8E55BD}
\newcommand{\gatonaranjo}[1][0.07]{\gato[#1]{gNar}{gNarS}}
\newcommand{\gatogris}[1][0.07]{\gato[#1]{gGri}{gGriS}}
\newcommand{\gatoblanco}[1][0.07]{\gato[#1]{gBla}{gBlaS}}
\newcommand{\gatonegro}[1][0.07]{\gato[#1]{gNeg}{gNegS}}
\newcommand{\gatoatigrado}[1][0.07]{\gato[#1]{gAti}{gAtiS}}
\newcommand{\gatolila}[1][0.07]{\gato[#1]{gLil}{gLilS}}

%% \paginagatos{color}{colorsuave}: última página del libro, los seis gatos con un globo «¡Tú puedes!» cada uno.
%% #1 color principal del libro (nar o ciencia), #2 su fondo suave (hielo o cienciaT). En blanco y negro sale en grises sola.
%% \globogato{color}{coordenada}{gato}: un gato centrado en la coordenada y su globo arriba a la derecha, con la colita apuntando al gato.
\newcommand{\globogato}[3]{%
  \node (g) at (#2) {#3[0.16]};
  \node[draw=#1,line width=1pt,fill=white,rounded corners=10pt,inner xsep=12pt,inner ysep=8pt,font=\Large\bfseries,text=#1,anchor=south west](b) at ([xshift=1.4cm,yshift=1.1cm]g.center) {¡Tú puedes!};
  \fill[white] ([xshift=14pt,yshift=0.6pt]b.south west) -- ++(-8pt,-9pt) -- ++(16pt,0) -- cycle;
  \draw[#1,line width=1pt] ([xshift=14pt]b.south west) -- ++(-8pt,-9pt) -- ++(16pt,9pt);}
\newcommand{\paginagatos}[2]{\clearpage\thispagestyle{empty}%
  \begin{tikzpicture}[remember picture,overlay]
    \fill[#2] (current page.north west) rectangle (current page.south east);
    \ifbn\draw[black,line width=2pt] ([xshift=1.5cm,yshift=-1.5cm]current page.north west) rectangle ([xshift=-1.5cm,yshift=1.5cm]current page.south east);\fi
    \node[anchor=north,text=#1,font=\fontsize{34}{38}\selectfont\bfseries] at ([yshift=-2.6cm]current page.north) {Un mensaje de los gatitos};
    \node[anchor=north,text=tinta,font=\Large,text width=15cm,align=center] at ([yshift=-4.4cm]current page.north) {Llegaste hasta el final del libro. Ellos no tienen dudas de lo que viene:};
    \globogato{#1}{[xshift=3.4cm,yshift=-8.6cm]current page.north west}{\gatonaranjo}
    \globogato{#1}{[xshift=12.2cm,yshift=-8.6cm]current page.north west}{\gatogris}
    \globogato{#1}{[xshift=3.4cm,yshift=-14.6cm]current page.north west}{\gatoblanco}
    \globogato{#1}{[xshift=12.2cm,yshift=-14.6cm]current page.north west}{\gatonegro}
    \globogato{#1}{[xshift=3.4cm,yshift=-20.6cm]current page.north west}{\gatoatigrado}
    \globogato{#1}{[xshift=12.2cm,yshift=-20.6cm]current page.north west}{\gatolila}
    \node[anchor=south,text=gris,font=\small] at ([yshift=1.8cm]current page.south) {Fabián Ramírez · Fabimath · fabimath.github.io/FabimathOficialWebSite/paes.html};
  \end{tikzpicture}\null\clearpage}
